% GENERATED FILE. Edit canonical lessons, then run npm run generate:books.
% canonical-source-hash: fnv1a64:949b8c3ee396276c
% canonical-lessons: ZH-C379-loudao, ZH-C379-anxin, ZH-W379-an-calm, ZH-C379-buan, ZH-C379-menjin, ZH-C379-zhian, ZH-R379-review-374-379

\chapter{Feeling Safe}
\label{ch:zh-feeling-safe}
\hlchaptermodality{379}

\begin{chapteropening}
\textbf{By the end of this chapter:} \emph{I can talk about the stairwell, feeling at ease or uneasy, the door system and law and order.}

Complete the last lesson of chapter 379: I can talk about the stairwell, feeling at ease or uneasy, the door system and law and order.
\end{chapteropening}

\section[lóudào]{lóudào — stairwell; corridor (in a block of flats)}

\label{lesson:ZH-C379-loudao}



\begin{quote}
Before the new one: say the Mandarin for a stranger, then the Mandarin for the property management office.
\end{quote}

\subsection*{You'll want to know: lóudào}
\emph{lóudào} — \textquotedblleft{}stairwell; corridor (in a block of flats)\textquotedblright{}.

\begin{quote}
\emph{lóudào lǐ de dēng huài le} — The light in the stairwell is broken.
\end{quote}

Its characters come later; for now it is a word for the ear.

\begin{cousinweb}
\emph{lóudào} is built from \emph{lóu} 'a building of several floors' and \emph{dào} 'a way', both inherited from Middle Chinese, \emph{luw} and \emph{dawX} (Guangyun, in Baxter's transcription), which glosses the first 'a house on a house' and the second 'a road'.
\end{cousinweb}

\begin{grammarlens}[title={What you've built}]
One more word for giving reasons.
\end{grammarlens}

\subsection*{Your turn}
\begin{itemize}
\raggedright
  \item \emph{Say it:} \emph{lóudào}
  \item \emph{Say it:} \emph{lóudào}, once more
  \item \emph{Say it:} \emph{lóudào lǐ de dēng huài le} — The light in the stairwell is broken.
  \item \emph{Recall:} say the Mandarin for a stranger, then the Mandarin for the property management office, then say \emph{lóudào} again
\end{itemize}

\begin{tcolorbox}[breakable,colback=teal!4,colframe=teal!35!black,title={Before you move on}]
What does lóudào mean? (\textquotedblleft{}Stairwell; corridor (in a block of flats)\textquotedblright{}.) Say it once more.
\end{tcolorbox}
\section[\texorpdfstring{\zh{安心} (ānxīn)}{安心 (ānxīn)}]{\zh{安心} (ānxīn) — with a mind at rest; settled}

\label{lesson:ZH-C379-anxin}



\begin{quote}
Before the new one: say the Mandarin for the property management office, then the Mandarin for the stairwell of a block.
\end{quote}

\subsection*{You'll want to know: \zh{安心}}
\textbf{\zh{安心}} — \emph{ānxīn} — \textquotedblleft{}with a mind at rest; settled\textquotedblright{}.

\begin{quote}
\textbf{\zh{这个社区很好我很安心}} \emph{zhège shèqū hěn hǎo wǒ hěn ānxīn} — This is a good neighbourhood, and my mind is at rest.
\end{quote}

The new character, in pieces: \zh{安} — six strokes: the roof of \zh{家} on top, three strokes, and \zh{女} under it. The next lesson puts it in your hand.

\begin{cousinweb}
\textbf{\zh{安心}} (\emph{ānxīn}) is built from \zh{安} 'calm' and \zh{心} 'the heart', both inherited from Middle Chinese, \emph{'an} and \emph{sim} (Guangyun, in Baxter's transcription), which glosses the first 'peaceful' and 'to come to rest'.
\end{cousinweb}

\begin{grammarlens}[title={What you've built}]
One more word for giving reasons.
\end{grammarlens}

\subsection*{Your turn}
\begin{itemize}
\raggedright
  \item \emph{Say it:} \emph{ānxīn}
  \item \emph{Say it:} \emph{ānxīn}, once more
  \item \emph{Say it:} \emph{zhège shèqū hěn hǎo wǒ hěn ānxīn} — This is a good neighbourhood, and my mind is at rest.
  \item \emph{Recall:} say the Mandarin for the property management office, then the Mandarin for the stairwell of a block, then say \emph{ānxīn} again
\end{itemize}

\begin{tcolorbox}[breakable,colback=teal!4,colframe=teal!35!black,title={Before you move on}]
What does \zh{安心} mean? (\textquotedblleft{}With a mind at rest; settled\textquotedblright{}.) Say it once more.
\end{tcolorbox}
\section[\texorpdfstring{\zh{安} (ān)}{安 (ān)}]{\zh{安} — the roof of \zh{家}, then \zh{女}}

\label{lesson:ZH-W379-an-calm}



\begin{quote}
The word first, then the character inside it.

\begin{itemize}
\raggedright
  \item \emph{Recall:} say \emph{ānxīn} — \textbf{R1}, one lesson back, before it decays
\end{itemize}

\begin{quote}
\zh{家} \zh{女}
\end{quote}

\emph{Read it:} each one aloud, with nothing said first — \emph{jiā}, \emph{nǚ}
\end{quote}

\subsection*{Script — the roof of \zh{家}, then \zh{女}}
\hlblockfigure{\detokenize{figures/ZH-W379-an-calm-filmstrip.pdf}}{How \zh{安} is written, stroke by stroke}

\begin{quote}
\zh{安}
\end{quote}

\textbf{Two pieces, top before bottom.}

\textbf{The roof, three strokes,} as on top of \zh{家} and \zh{字}: a dot at the top, a short stroke down on the left, then a bar across that turns down at the right end and hooks.

\textbf{Under the roof, \zh{女}, three strokes,} as you wrote it: a stroke that falls to the left and turns to the right, a long stroke falling to the left, then a long bar across.

3 + 3 = six.

\subsection*{Your turn}
\begin{itemize}
\raggedright
  \item \emph{Write it:} \zh{安}, all six, with the model in front of you
  \item \emph{Cover:} write it again from memory, and count the pieces, not the strokes
\end{itemize}

\begin{tcolorbox}[breakable,colback=teal!4,colframe=teal!35!black,title={Before you move on}]
How many strokes in \zh{安}? (\textbf{Six}.) What is on top? (\textbf{The roof of \zh{家}}.) What sits under it? (\textbf{\zh{女}}.)

Source: \href{https://raw.githubusercontent.com/chanind/hanzi-writer-data/68d10a4b21150cae5e1ebbd223eed289cf32d90c/data/\%E5\%AE\%89.json}{Hanzi Writer Data \zh{安}, pinned stroke paths and medians}.
\end{tcolorbox}
\section[\texorpdfstring{\zh{不安} (bù'ān)}{不安 (bù'ān)}]{\zh{不安} (bù'ān) — uneasy; worried}

\label{lesson:ZH-C379-buan}



\begin{quote}
Before the new one: say the Mandarin for the stairwell of a block, then the Mandarin for with a mind at rest.
\end{quote}

\subsection*{You'll want to know: \zh{不安}}
\textbf{\zh{不安}} — \emph{bù'ān} — \textquotedblleft{}uneasy; worried\textquotedblright{}.

\begin{quote}
\textbf{\zh{我一个人在家的时候有点儿不安}} \emph{wǒ yí ge rén zài jiā de shíhou yǒudiǎnr bù'ān} — When I'm at home on my own, I feel a little uneasy.
\end{quote}

Tone: \zh{不} keeps its own falling tone here, bù, because ān after it is a level tone, not a falling one.

\begin{cousinweb}
\textbf{\zh{不安}} (\emph{bù'ān}) is built from \zh{不} 'not' and \zh{安} 'calm', both inherited from Middle Chinese, \emph{pjut} and \emph{'an} (Guangyun, in Baxter's transcription).
\end{cousinweb}

\begin{grammarlens}[title={What you've built}]
One more word for giving reasons.
\end{grammarlens}

\subsection*{Your turn}
\begin{itemize}
\raggedright
  \item \emph{Say it:} \emph{bù'ān}
  \item \emph{Say it:} \emph{bù'ān}, once more
  \item \emph{Say it:} \emph{wǒ yí ge rén zài jiā de shíhou yǒudiǎnr bù'ān} — When I'm at home on my own, I feel a little uneasy.
  \item \emph{Recall:} say the Mandarin for the stairwell of a block, then the Mandarin for with a mind at rest, then say \emph{bù'ān} again
\end{itemize}

\begin{tcolorbox}[breakable,colback=teal!4,colframe=teal!35!black,title={Before you move on}]
What does \zh{不安} mean? (\textquotedblleft{}Uneasy; worried\textquotedblright{}.) Say it once more.
\end{tcolorbox}
\section[ménjìn]{ménjìn — door entry system; a card-entry door}

\label{lesson:ZH-C379-menjin}



\begin{quote}
Before the new one: say the Mandarin for with a mind at rest, then the Mandarin for uneasy (ill at ease).
\end{quote}

\subsection*{You'll want to know: ménjìn}
\emph{ménjìn} — \textquotedblleft{}door entry system; a card-entry door\textquotedblright{}.

\begin{quote}
\emph{wǒmen xiǎoqū yǒu ménjìn suǒyǐ hěn ānquán} — Our estate has a door entry system, so it's safe.
\end{quote}

Its characters come later; for now it is a word for the ear.

\begin{cousinweb}
\emph{ménjìn} is built from \zh{门} 'a gate' and \emph{jìn} 'to forbid', both inherited from Middle Chinese, \emph{mwon} and \emph{kimH} (Guangyun, in Baxter's transcription), which glosses the second 'to restrain' and 'to stop'. The card-and-lock system is a modern use.
\end{cousinweb}

\begin{grammarlens}[title={What you've built}]
One more word for giving reasons.
\end{grammarlens}

\subsection*{Your turn}
\begin{itemize}
\raggedright
  \item \emph{Say it:} \emph{ménjìn}
  \item \emph{Say it:} \emph{ménjìn}, once more
  \item \emph{Say it:} \emph{wǒmen xiǎoqū yǒu ménjìn suǒyǐ hěn ānquán} — Our estate has a door entry system, so it's safe.
  \item \emph{Recall:} say the Mandarin for with a mind at rest, then the Mandarin for uneasy (ill at ease), then say \emph{ménjìn} again
\end{itemize}

\begin{tcolorbox}[breakable,colback=teal!4,colframe=teal!35!black,title={Before you move on}]
What does ménjìn mean? (\textquotedblleft{}Door entry system; a card-entry door\textquotedblright{}.) Say it once more.
\end{tcolorbox}
\section[zhì'ān]{zhì'ān — law and order; public order}

\label{lesson:ZH-C379-zhian}



\begin{quote}
Before the new one: say the Mandarin for uneasy (ill at ease), then the Mandarin for a card-entry door.
\end{quote}

\subsection*{You'll want to know: zhì'ān}
\emph{zhì'ān} — \textquotedblleft{}law and order; public order\textquotedblright{}.

\begin{quote}
\emph{zhèr de zhì'ān hěn hǎo suǒyǐ wǒ hěn fàngxīn} — Law and order is good here, so I don't worry.
\end{quote}

Its characters come later; for now it is a word for the ear.

\begin{cousinweb}
\emph{zhì'ān} is built from \emph{zhì} 'to govern' and \zh{安} 'calm', both inherited from Middle Chinese, \emph{dri} and \emph{'an} (Guangyun, in Baxter's transcription), which glosses the first 'to put in order' and also lists it twice in the departing tone, \emph{drijH} and \emph{driH}, with the same gloss. Mandarin zhì follows a departing reading.
\end{cousinweb}

\begin{grammarlens}[title={What you've built}]
One more word for giving reasons.
\end{grammarlens}

\subsection*{Your turn}
\begin{itemize}
\raggedright
  \item \emph{Say it:} \emph{zhì'ān}
  \item \emph{Say it:} \emph{zhì'ān}, once more
  \item \emph{Say it:} \emph{zhèr de zhì'ān hěn hǎo suǒyǐ wǒ hěn fàngxīn} — Law and order is good here, so I don't worry.
  \item \emph{Recall:} say the Mandarin for uneasy (ill at ease), then the Mandarin for a card-entry door, then say \emph{zhì'ān} again
\end{itemize}

\begin{tcolorbox}[breakable,colback=teal!4,colframe=teal!35!black,title={Before you move on}]
What does zhì'ān mean? (\textquotedblleft{}Law and order; public order\textquotedblright{}.) Say it once more.
\end{tcolorbox}
\section[(dialogue)]{Looking back: from People and Government to Feeling Safe}

\label{lesson:ZH-R379-review-374-379}



\begin{quote}
Say the Mandarin for the population (of a place) and the government.
\end{quote}

\subsection*{Your turn}
\begin{itemize}
\raggedright
  \item \emph{Say it:} the population (of a place), the government, a people's culture, an ethnic group, to develop and grow
  \item \emph{Say it:} poor (having little money), society, to donate (money), a volunteer, each other
  \item \emph{Say it:} an organised event, a social worker, to treat with care, to give up one's seat, square dancing (in a public square)
  \item \emph{Say it:} a resident (of a place), a local community, a housing estate, the area around (a place), a home near a good school
  \item \emph{Say it:} unwanted noise, to grumble about, soundproofing (keeping sound out), a stranger, the property management office
  \item \emph{Say it:} the stairwell of a block, with a mind at rest, uneasy (ill at ease), a card-entry door, law and order
  \item \emph{Read it:} \zh{社} \zh{区} \zh{安} — each one aloud, then a word that holds it
\end{itemize}

\begin{tcolorbox}[breakable,colback=teal!4,colframe=teal!35!black,title={Before you move on}]
The population (of a place)? (\textbf{\zh{人口}}, \emph{rénkǒu}.) The government? (\emph{zhèngfǔ}.) A people's culture? (\emph{wénhuà}.) An ethnic group? (\emph{mínzú}.) To develop and grow? (\emph{fāzhǎn}.) Poor (having little money)? (\emph{qióng}.) Society? (\textbf{\zh{社会}}, \emph{shèhuì}.) To donate (money)? (\emph{juān}.) A volunteer? (\emph{zhìyuànzhě}.) Each other? (\emph{hùxiāng}.) An organised event? (\emph{huódòng}.) A social worker? (\textbf{\zh{社工}}, \emph{shègōng}.) To treat with care? (\emph{àihù}.) To give up one's seat? (\emph{ràngzuò}.) Square dancing (in a public square)? (\emph{guǎngchǎngwǔ}.) A resident (of a place)? (\emph{jūmín}.) A local community? (\textbf{\zh{社区}}, \emph{shèqū}.) A housing estate? (\textbf{\zh{小区}}, \emph{xiǎoqū}.) The area around (a place)? (\emph{zhōuwéi}.) A home near a good school? (\textbf{\zh{学区房}}, \emph{xuéqūfáng}.) Unwanted noise? (\emph{zàoyīn}.) To grumble about? (\emph{bàoyuàn}.) Soundproofing (keeping sound out)? (\emph{géyīn}.) A stranger? (\textbf{\zh{生人}}, \emph{shēngrén}.) The property management office? (\emph{wùyè}.) The stairwell of a block? (\emph{lóudào}.) With a mind at rest? (\textbf{\zh{安心}}, \emph{ānxīn}.) Uneasy (ill at ease)? (\textbf{\zh{不安}}, \emph{bù'ān}.) A card-entry door? (\emph{ménjìn}.) Law and order? (\emph{zhì'ān}.)
\end{tcolorbox}
